\documentclass[10pt,a4paper]{amsart}
\usepackage[margin=19mm]{geometry}
\usepackage{amsmath,amssymb,mathtools}
\usepackage[scr=rsfs,cal=euler]{mathalfa}
\usepackage{xfrac}
\usepackage[unicode,hidelinks,bookmarks=false]{hyperref}
\newcommand{\ie}{i.e.\ }
\newcommand{\cf}{cf.\ }
\newcommand{\resp}{respectively\ }
\newcommand{\Subsection}{\subsection}
\providecommand{\nobreakdash}{\nobreak\hbox{-}}
\setlength{\emergencystretch}{2em}
\multlinegap0pt
\usepackage{amssymb}
\usepackage{tikz}
\newtheorem{lem}{Lemma}
\title{Model theory and Connes' bicentralizer problem}
\author{Hiroshi Ando, Isaac Goldbring}
\date{Source: arXiv:2605.12776v1 (CC BY 4.0)}
\begin{document}
\maketitle

\noindent\emph{Source and licence.} Model theory and Connes' bicentralizer problem. Version arXiv:2605.12776v1 (CC BY 4.0), distributed by arXiv under the Creative Commons Attribution 4.0 International licence, http://creativecommons.org/licenses/by/4.0/, as recorded in the arXiv licence metadata for this exact version and shown on the abstract page https://arxiv.org/abs/2605.12776v1. Full licence notice: \texttt{licenses/arxiv-2605.12776-CC-BY-4.0.txt}.\par\smallskip

\noindent\textit{Editorial scope.} Counted unit: the article's lem supp (source0.tex lines 1501--1511). The article's complete original proof of it is delivered with it (source0.tex lines 1812--2012, one \texttt{proof} environment of 201 lines, which the article places away from the statement). No mathematics is written, repaired or re-derived: the statement and the proof are the article's own lines, in the article's own notation, and no retained line is substituted or re-flowed.  Retained uncounted as the source-local context the proof needs: lines 1698--1718.  Nothing was omitted from any retained span, so the package carries no omission record.  No retained line contains a citation, so the wrapper carries no bibliography and no bibliography entry.  No retained line contains a figure environment, an image-inclusion command or drawing code, so no image is delivered and the package declares no asset dependency.  Editorial formatting only, in addition: an AMS \texttt{amsart} preamble that restates the article's own declarations and macros used by the retained lines, namely the environments \texttt{lem} on separate counters, together with the standard packages those lines need.  The retained environments print in this document as Lemma 1 in order of appearance, whereas the article prints the same items with the numbering of its own full document.  The complete native source of the article as distributed in the arXiv e-print for this exact version is bundled unchanged as \texttt{sources/200450/source0.tex}.
\begin{lem}\label{lem one-sided spec estimate}
Let \(h\in C_b(\mathbb R)\), and regard \(h\) as a tempered distribution. Assume
that
\[
        \operatorname{supp}(\widehat h)\subset [0,b]
\]
for some \(b\ge 0\). Then \(h\) has an entire extension \(H\), and
\[
        |H(z)|\le \|h\|_\infty,
        \qquad \operatorname{Im}z\le 0.
\]
Similarly, if
\[
        \operatorname{supp}(\widehat h)\subset [-b,0],
\]
then \(h\) has an entire extension \(H\), and
\[
        |H(z)|\le \|h\|_\infty,
        \qquad \operatorname{Im}z\ge 0.
\]
\end{lem}

\begin{lem}\label{lem FT cpt supp}
Let $f\in C_b(\mathbb R)$ and regard $f$ as a tempered distribution. Assume that
\[
\operatorname{supp}(\widehat{f})\subset [-a,a]
\]
for some $a\ge 0$. Then $f$ extends to an entire function on $\mathbb C$, still
denoted by $f$, and for every $s\in \mathbb{R}$, we have
\[
\sup_{t\in\mathbb R}|f(t-is)|\le e^{a|s|}\|f\|_\infty.
\]
\end{lem}

\begin{proof}[Proof of Lemma \ref{lem FT cpt supp}]
First consider the lower half-plane direction. Define
\[
        h(t):=e^{-iat}f(t).
\]
Then
\[
        h\in C_b(\mathbb R),
        \qquad
        \|h\|_\infty=\|f\|_\infty.
\]
With the Fourier transform convention
\[
        \widehat u(\lambda)
        =
        \int_{\mathbb R}e^{it\lambda}u(t)\,dt,
\]
we see that in \(\mathcal S'(\mathbb R)\),
\[
        \widehat h(\lambda)
        =
        \widehat f(\lambda-a).
\]

By definition of equality of tempered distributions, this means that for every
\(\varphi\in\mathcal S(\mathbb R)\),
\[
        \langle \widehat h,\varphi\rangle
        =
        \langle \widehat f,\varphi(\cdot+a)\rangle .
\]

Indeed, recall that the Fourier transform of a tempered distribution is defined by
\[
        \langle \widehat T,\varphi\rangle
        =
        \langle T,\widehat\varphi\rangle,
        \qquad \varphi\in\mathcal S(\mathbb R).
\]
Therefore
\[
\begin{aligned}
        \langle \widehat h,\varphi\rangle
        &=
        \langle h,\widehat\varphi\rangle                                      \\
        &=
        \left\langle e^{-iat}f(t),
        \int_{\mathbb R}e^{it\lambda}\varphi(\lambda)\,d\lambda
        \right\rangle                                                         \\
        &=
        \left\langle f(t),
        \int_{\mathbb R}e^{it(\lambda-a)}\varphi(\lambda)\,d\lambda
        \right\rangle .
\end{aligned}
\]
Putting
\[
        \mu=\lambda-a,
        \qquad \lambda=\mu+a,
\]
we get
\[
        \int_{\mathbb R}e^{it(\lambda-a)}\varphi(\lambda)\,d\lambda
        =
        \int_{\mathbb R}e^{it\mu}\varphi(\mu+a)\,d\mu.
\]
Thus
\[
\begin{aligned}
        \langle \widehat h,\varphi\rangle
        &=
        \left\langle f(t),
        \int_{\mathbb R}e^{it\mu}\varphi(\mu+a)\,d\mu
        \right\rangle                                                         \\
        &=
        \langle \widehat f,\varphi(\cdot+a)\rangle .
\end{aligned}
\]
This proves
\[
        \widehat h(\lambda)=\widehat f(\lambda-a)
\]
in \(\mathcal S'(\mathbb R)\).
Therefore
\[
        \operatorname{supp}(\widehat h)
        \subset [0,2a].
\]
By Lemma \ref{lem one-sided spec estimate}, \(h\) has an entire extension \(H\) such
that
\[
        |H(z)|\le \|h\|_\infty=\|f\|_\infty,
        \qquad \operatorname{Im}z\le 0.
\]
Define
\[
        F_-(z):=e^{iaz}H(z).
\]
Then \(F_-\) is entire and, for real \(t\),
\[
        F_-(t)
        =
        e^{iat}H(t)
        =
        e^{iat}h(t)
        =
        f(t).
\]
Thus \(F_-\) is an entire extension of \(f\). For \(s\ge 0\), we have
\[
        |F_-(t-is)|
        =
        |e^{ia(t-is)}|\,|H(t-is)|
        =
        e^{as}|H(t-is)|
        \le
        e^{as}\|f\|_\infty .
\]
Hence
\[
        \sup_{t\in\mathbb R}|F_-(t-is)|
        \le
        e^{as}\|f\|_\infty,
        \qquad s\ge 0.
\]

Now consider the upper half-plane direction. Define
\[
        k(t):=e^{iat}f(t).
\]
Then
\[
        k\in C_b(\mathbb R),
        \qquad
        \|k\|_\infty=\|f\|_\infty,
\]
and
\[
        \widehat k(\lambda)
        =
        \widehat f(\lambda+a).
\]
Therefore
\[
        \operatorname{supp}(\widehat k)
        \subset [-2a,0].
\]
By Lemma \ref{lem one-sided spec estimate} again, \(k\) has an entire extension \(K\) such
that
\[
        |K(z)|\le \|k\|_\infty=\|f\|_\infty,
        \qquad \operatorname{Im}z\ge 0.
\]
Define
\[
        F_+(z):=e^{-iaz}K(z).
\]
Then \(F_+\) is entire and, for real \(t\),
\[
        F_+(t)
        =
        e^{-iat}K(t)
        =
        e^{-iat}k(t)
        =
        f(t).
\]
Thus \(F_+\) is also an entire extension of \(f\). Since \(F_-\) and \(F_+\)
agree on \(\mathbb R\), they agree on all of \(\mathbb C\) by the identity
theorem. Denote this common entire extension simply by \(f\).

If \(s<0\), put \(r=-s>0\). Then \(t-is=t+ir\), and
\[
        |f(t-is)|
        =
        |f(t+ir)|
        =
        |F_+(t+ir)|
        =
        |e^{-ia(t+ir)}|\,|K(t+ir)|
        =
        e^{ar}|K(t+ir)|
        \le
        e^{ar}\|f\|_\infty.
\]
Since \(r=|s|\), this gives
\[
        |f(t-is)|
        \le
        e^{a|s|}\|f\|_\infty
        \qquad (s<0).
\]

Combining the estimates for \(s\ge 0\) and \(s<0\), we obtain
\[
        \sup_{t\in\mathbb R}|f(t-is)|
        \le
        e^{a|s|}\|f\|_\infty
\]
for every \(s\in\mathbb R\).
\end{proof}

\end{document}
